% =====================================================================
% 8. Dynamics of VaR and ES -- question (ix)
% =====================================================================
\section{Dynamics of VaR and ES (Question ix)}
\label{sec:rolling}

\subsection{DESIGN}
At each estimation date we take the last $500$ daily log-returns, about two years, and re-estimate the three models. A window of $500$ days contains only $100$ non-overlapping weeks, too few for a $t$-copula and for an empirical $99.9\%$ quantile. Within each window we therefore use the $496$ \emph{overlapping} five-day log-returns. The estimators remain consistent, but the observations are serially dependent, so the likelihood is a pseudo-likelihood and standard errors are not reported. The window is moved forward by five trading days, which gives \val{q9.n_windows|int} estimation dates from \val{q9.first|date} to \val{q9.last|date}. At every date $s_k$ and the reference thresholds $d_k$ are recomputed from the re-estimated model with the default probabilities $\pi_k$ unchanged (closed form under M2, exact average under M1, two-dimensional quadrature with step $0.2$ under M3), and the loss distribution is simulated with $50{,}000$ scenarios. The reference thresholds are used because the purpose is to compare dates and models, and they keep every default probability at exactly $\pi_k$ in every window. The same random numbers are used at all dates and for all models, so that changes through time are not caused by a new draw of the simulation noise.

To read the differences between M3 and M2 we need to know how large they would be if the two models were identical. We therefore run a \emph{placebo experiment} (Appendix~B): M3 is given the parameters of M2 and $\nu=200$, at which the $t$-copula is Gaussian for practical purposes, so that the true difference is zero; the difference in VaR and ES is then measured exactly as in one rolling window, for $200$ independent seeds. The relative difference in VaR has mean \val{q9.placebo.VaR_mean|spct2}\% and standard deviation \val{q9.placebo.VaR_sd|pct1}\%; for ES the figures are \val{q9.placebo.ES_mean|spct2}\% and \val{q9.placebo.ES_sd|pct1}\%. The standard deviation is the simulation noise against which a rolling difference has to be judged.

\subsection{RESULTS}
\begin{table}[tb]
\caption{Rolling $99.9\%$ VaR and ES (USD), \val{q9.n_windows|int} windows of 500 daily returns, thresholds recalibrated in every window.}
\label{tab:rolling}
\centering
\qrmtable{tab_rolling}
\end{table}

\begin{figure}[p]
\centering
\qrmfigure{fig_rolling}
\caption{Rolling estimates on windows of 500 daily returns. Top two panels: $99.9\%$ VaR and ES with thresholds recalibrated to $\pi_k$ in every window. Third panel: VaR when the thresholds are frozen at their full-sample values. Lower panels: dependence parameters and degrees of freedom of the $t$-copula.}
\label{fig:rolling}
\end{figure}

\begin{figure}[tb]
\centering
\qrmfigure{fig_rolling_diff}
\caption{Relative difference between M3 and M2 in the rolling $99.9\%$ VaR (top) and ES (bottom). Grey band: mean $\pm$ two standard deviations of the same difference in the placebo experiment, in which the two models coincide.}
\label{fig:rollingdiff}
\end{figure}

Figure~\ref{fig:rolling} and Table~\ref{tab:rolling} show the results. The three models behave very differently through time.

\paragraph{M2 is almost constant.}
The Gaussian VaR has a mean of USD~\val{q9.mean.VaR_M2|usd} and a standard deviation across windows of only USD~\val{q9.std.VaR_M2|usd}, less than 1\%; it stays between USD~\val{q9.min.VaR_M2|usd} and USD~\val{q9.max.VaR_M2|usd} although the weekly volatility of the indices varies by a factor of \val{q9.sigma_ratio_max_min|f1} across windows (from \val{q9.min.sigma_SPX|pct1}\% to \val{q9.max.sigma_SPX|pct1}\% for SPX). This is the invariance property (\ref{eq:standardised}) at work. Under M2 the standardised combinations $V_k$ are exactly Gaussian in every window, so once the thresholds are recalibrated to the same $\pi_k$ the data enter only through the correlation between the two indices and the ratio of their volatilities. The estimated correlation moves between \val{q9.min.rho_gauss|f2} and \val{q9.max.rho_gauss|f2}, and the correlation between the rolling VaR and the rolling $\hat\rho$ is \val{q9.corr_VaRM2_rho|f2}: a higher correlation between the factors means less diversification between Swiss-loaded and US-loaded names.

\paragraph{M3 against M2: inside the simulation noise for VaR, slightly above it for ES.}
Figure~\ref{fig:rollingdiff} plots the relative difference between M3 and M2 against the placebo band. For VaR the difference lies between \val{q9.dVaR_min|spct1}\% and \val{q9.dVaR_max|spct1}\%, averages \val{q9.dVaR_mean|spct1}\%, and is negative in \val{q9.share_VaR_M3_below_M2|pct0}\% of the windows; \val{q9.share_dVaR_in_band|pct0}\% of the windows lie inside the band. In the rolling analysis the VaR of M3 is therefore not systematically above that of M2, and the difference in a single window cannot be interpreted. Two features of the design explain the pattern. First, the degrees of freedom are poorly identified in two-year windows: $\hat\nu$ is at its upper bound of 200 in \val{q9.share_nu_at_bound|pct0}\% of the windows and below 10 in \val{q9.share_nu_below_10|pct0}\% of them, typically those that contain a crisis. At the bound M3 is M2 for practical purposes, yet the measured VaR difference averages \val{q9.dVaR_mean_bound|spct1}\% with a standard deviation of only \val{q9.dVaR_sd_bound|pct1}\%. The reason is that the same random numbers are used at every date: the simulation noise is not redrawn from window to window but is, for given parameters, a fixed offset, which for our seed is negative and well inside the placebo band. Relative to that offset, the windows with $\hat\nu<10$ show a VaR difference that is \val{q9.dVaR_low_minus_bound|pct1} percentage points higher, comparable in size to the full-sample effect of \val{q8.M3.VaR_vs_M2|pct1}\%. Second, no result guarantees the sign of the difference once the thresholds are recalibrated: the joint estimation of M3 also changes $\hat\rho$ and the margins relative to M2, and the correlation between the VaR difference and $\hat\rho_{t}-\hat\rho_{\mathrm{Gauss}}$ across windows is \val{q9.corr_dVaR_drho|f2}.

For ES the picture is clearer. The difference averages \val{q9.dES_mean|spct1}\%, reaches \val{q9.dES_max|spct1}\%, is negative in only \val{q9.share_ES_M3_below_M2|pct0}\% of the windows and lies above the placebo band in \val{q9.share_dES_above_band|pct0}\% of them; in the windows with $\hat\nu<10$ it averages \val{q9.dES_mean_nu_below_10|spct1}\%, and its correlation with the tail-dependence coefficient $\lambda_L$ of (\ref{eq:taildep}) is \val{q9.corr_dES_taildep|f2}. Tail dependence moves probability mass beyond the $99.9\%$ quantile more than it moves the quantile itself, so it shows in ES before it shows in VaR. Even so, with $\hat\nu$ around 5 and $\lambda_L$ above $0.5$ the ES of M3 exceeds that of M2 by less than 10\%. The rolling analysis thus confirms the full-sample finding: the copula of two Gaussian factors has little leverage on the tail of the loss.

\paragraph{M1 depends on which weeks are in the window.}
The empirical VaR has a mean of USD~\val{q9.mean.VaR_M1|usd} and a standard deviation of USD~\val{q9.std.VaR_M1|usd}, \val{q9.ratio_std_M1_M2|f0} times that of M2. It ranges from USD~\val{q9.min.VaR_M1|usd} (window ending \val{q9.argmin_VaR_M1|date}) to USD~\val{q9.max.VaR_M1|usd} (window ending \val{q9.argmax_VaR_M1|date}), that is from \val{q9.ratio_M1_M2_min|ratio_pct0}\% to \val{q9.ratio_M1_M2_max|ratio_pct0}\% of the VaR under M2. With $496$ equally likely observations each has probability $0.2\%$, so the $99.9\%$ tail of the factor distribution is the single worst five-day return of the window. The VaR therefore jumps when a crash enters the window (autumn 2008, August 2011, March 2020) and falls back abruptly 500 days later when it drops out, the familiar ghost effect of historical simulation. Its correlation with the average excess kurtosis of the two indices in the window is \val{q9.corr_VaRM1_kurt|f2}. In calm windows without any crash the empirical distribution can even be thinner-tailed than the Gaussian: M1 is below M2 in \val{q9.share_M1_below_M2|pct0}\% of the windows. At the latest date, \val{q9.last|date}, the VaR is USD~\val{q9.last_row.VaR_M1|usd} under M1, USD~\val{q9.last_row.VaR_M2|usd} under M2 and USD~\val{q9.last_row.VaR_M3|usd} under M3.

\paragraph{What recalibration hides.}
The stability of M2 and M3 is a consequence of the design and not a property of credit risk. Because the thresholds are reset in every window so that default probabilities equal the given $\pi_k$, a rise in market volatility is absorbed entirely by lower thresholds and leaves the loss distribution unchanged. The model is through-the-cycle by construction. To show what this removes, the third panel of Figure~\ref{fig:rolling} reports VaR when the thresholds are frozen at their full-sample values, so that a window with more volatile or lower returns implies higher default probabilities. In this point-in-time reading the Gaussian VaR ranges from USD~\val{q9.min.VaRfix_M2|usd} (window ending \val{q9.argmin_VaRfix_M2|date}) to USD~\val{q9.max.VaRfix_M2|usd} (\val{q9.argmax_VaRfix_M2|date}), its correlation with the level of volatility is \val{q9.corr_VaRfixM2_sigma|f2}, and the implied expected loss moves between USD~\val{q9.min.ELfix_M2|usd} and USD~\val{q9.max.ELfix_M2|usd}. The differences between the three models then become second-order relative to the common volatility cycle. Which reading is appropriate depends on what $\pi_k$ represents: a rating-based, through-the-cycle default probability calls for recalibration, as the assignment prescribes; a market-implied one would call for fixed barriers.
